% 20240922.typ（Typst）を LaTeX で再現したもの
% ratex v0.4.7 でのビルド: ratex -xelatex 20240922.tex
\documentclass[10pt,twocolumn]{article}

% 用紙：JIS B4 横，余白 左右 1.5cm・上下 1cm
\usepackage[paperwidth=364mm,paperheight=257mm,hmargin=1.5cm,vmargin=1cm]{geometry}
\setlength{\columnsep}{13mm}

\usepackage{amsmath,amssymb}
\usepackage{enumitem}
\usepackage{xeCJK}

% フォント：BIZ UDP明朝（本文），BIZ UDPゴシック（見出し・太字）
% ratex は OS のフォントを名前では探さないので，ファイルと面番号を指定する
% （BIZ UDP 系は .ttc の面 1）
\setmainfont{BIZ-UDMinchoM}[Path=/Library/Fonts/,Extension=.ttc,FontIndex=1]
\setsansfont{BIZ-UDGothicR}[Path=/Library/Fonts/,Extension=.ttc,FontIndex=1,BoldFont=BIZ-UDGothicB]
\setCJKmainfont{BIZ-UDMinchoM}[Path=/Library/Fonts/,Extension=.ttc,FontIndex=1]
\setCJKsansfont{BIZ-UDGothicR}[Path=/Library/Fonts/,Extension=.ttc,FontIndex=1,BoldFont=BIZ-UDGothicB]

% 段落：両端揃え，字下げなし，行送りを広めに
\setlength{\parindent}{0pt}
\setlength{\parskip}{2pt}
\AtBeginDocument{\fontsize{10pt}{17pt}\selectfont}
\pagestyle{empty}

% 小問の番号 (1), (2), ...
\setlist{leftmargin=*,labelindent=0pt,itemsep=0pt,parsep=0pt,topsep=0pt,partopsep=0pt}
\setlist[enumerate]{label=(\arabic*)}

% #dai：18pt ゴシック太字，太さ 1.5pt の下線
\newcommand{\dai}[1]{%
  {\fontsize{18pt}{26pt}\selectfont\sffamily\bfseries
   \sbox0{#1}\usebox0\hspace{-\wd0}\rule[-5.75pt]{\wd0}{1.5pt}}}

% #bf：ゴシック太字
\newcommand{\gbf}[1]{{\sffamily\bfseries #1}}

% #hako：幅 1cm・高さ 6mm の解答欄
\newcommand{\hako}[1]{{\setlength{\fboxsep}{0pt}%
  \fbox{\rule[-1.8mm]{0pt}{6mm}\makebox[10mm][l]{\hspace*{3.5mm}#1}}}}

% 見出し（= ...）：番号「1.」を大きめの太字ゴシックで。問題文は同じ行に続ける
\newcounter{mondai}
\newcommand{\mondai}{%
  \par\vspace{0.6em}\stepcounter{mondai}%
  \noindent{\fontsize{14pt}{17pt}\selectfont\sffamily\bfseries\themondai.}\ \ }

\begin{document}

\twocolumn[\dai{2024年度　第3学期　実力考査　中学3年　数学}\vspace{6mm}]

\noindent\hspace*{10cm}\gbf{2025年1月　　日（　）　担当：清水}

\gbf{【注意事項】}

\gbf{3. 〜 8. は答えのみではなく，考え方や計算式・図などを記すこと。}

\vspace{5mm}

\gbf{今回の実力試験は「距離」をテーマにして色々な問題について考えてみます。}

\vspace{5mm}

%--------------------------------------------------------------- 1
\mondai
\hako{　} を適切に埋めよ。（答えのみでよい。）

\vspace{5mm}

\gbf{清水：} 今回はいろいろな「距離」について考えてみましょう。まずは1次元で考えてみます。$x$ 軸上に2点A$(a)$ ，B$(b)$ があります。この2点の距離ABは絶対値を用いてどのように表されますか？

\gbf{太郎：} AB$=$ \hako{ア} となります。例えば，$a=10$，$b=-23$ のとき，AB$=$ \hako{イ} です。

\gbf{清水：} $x$ 軸上の3点をX$(x)$ ，A$(-1)$ ，B$(2)$ とします。AXの距離がBXの2倍であるときの実数 $x$ を求めてみましょう。

\gbf{花子：} 今回の関係式は
$\mathrm{AX}=2\mathrm{BX}$
となります。 $x$ についての式にすると，\hako{ウ} となり，この関係式を満たす $x$ は， $x=$ \hako{エ} です。$x$ は複数個あるのですね。

\vspace{5mm}

\gbf{清水：} 続いて2次元で考えてみましょう。まず，$xy$平面上の2点P$(a,b)$，Q$(c,d)$があるとき，2点P,Qの距離はどのように表すことができますか？

\gbf{花子：} PQ$=$ \hako{オ} となります。例えば，P$(-3,10)$，Q$(7,-7)$のとき，PQ$=$ \hako{カ} です。

\gbf{清水：} $xy$ 平面上の3点をX$(x,0)$ ，R$(1,4)$ ，S$(3,-6)$ とします。$\mathrm{RX}^2+\mathrm{SX}^2$の値が最小となる$x$のその時の値はどうなるでしょうか？

\gbf{太郎：} $x$を変量と考えるとと$\mathrm{RX}^2+\mathrm{SX}^2$は$x$の二次関数になりますね。$\mathrm{RX}^2+\mathrm{SX}^2$を$x$であらわすと， \hako{キ} となり，平方完成して， \hako{ク} となります。よって，$x=$ \hako{ケ} のとき最小値 \hako{コ} ですね。

\vspace{5mm}

\gbf{花子：} 清水先生，もし，$\mathrm{RX}+\mathrm{SX}$だったら，最小値はもとまるのでしょうか？

\gbf{清水：} 求められるよ。考えてみてください。

\gbf{花子：} $x$についての式で表してみたけど，今の私たちでは最小値が求められるような式ではないですね。

\gbf{太郎：} そうだね。とりあえず，図でも書いてみる？
\[
  \text{＜考え中＞}
\]

\gbf{太郎：} 図を書いて考えるとわかることもあるんだね。$x=$ \hako{サ} のとき，最小値は \hako{シ} とわかりました。

\gbf{清水：} よかったですね！

\newpage %-------------------------------------------------- 段の区切り

\gbf{距離を考えるときに制限がある場合があります。}

%--------------------------------------------------------------- 2
\mondai
図のような直方体ABCD-EFGHがあります。AB=3，AD = 4，AE=7です。辺AB上にAP=2となる点P，辺GC上にGQ=1となる点Qがあります。

\begin{enumerate}
  \item 直方体の辺上だけを移動できるとき，PからQへの最短距離を求めよ。
  \item 直方体の辺，面，及びその内部を移動できるとき，PからQへの最短距離を求めよ。
  \item 直方体の辺及び面上だけを移動できるとき，PからQへの最短距離を求めよ。
\end{enumerate}

\vspace{15mm}

%--------------------------------------------------------------- 3
\mondai
図のような格子状の経路があります。縦の長さは10，横の長さは6です。次の問いに答えよ。

\begin{enumerate}
  \item 点Aから点Bへの最短経路は何通りですか？
  \item 点Aから点Bへの最短経路のうち，PとQのどちらも通らない方法は何通りですか？
  \item 点Aから点Bへの最短経路のうち，4回曲がるのは何通りですか？
\end{enumerate}

\vspace{15mm}

\gbf{距離と角度について考えてみましょう}

%--------------------------------------------------------------- 4
\mondai
AB=10，AD=5の長方形があります。点Aから球を発車して，点Cへ到達到達します。次の(1)〜(3)の場合において，図の $\theta$ と球の移動した距離を求めよ。ただし，球は辺とぶつかると，入射角と反射角が同じになるように跳ね返ります。$\theta$ を求めるとき，三角関数表を用いてよい。

\begin{enumerate}
  \item \mbox{}
  \item \mbox{}
  \item \mbox{}
\end{enumerate}

\gbf{地球儀のように球面上で距離を考えることがあります。}

%--------------------------------------------------------------- 5
\mondai
半径が10cmの地球儀の赤道上の点Pにいます。Pからこの地球儀上を次のように移動します。

\begin{itemize}
  \item 点Pから北に10cm移動して点Qに到達
  \item 点Qから東に10cm移動して点Rに到達
  \item 点Rから南に10cm移動して点Sに到達
\end{itemize}

次の問いに答えよ。三角関数表を用いてよい。

\begin{enumerate}
  \item 点Pは北緯何度か？ただし，赤道は北緯0度，北極点は北緯90度である。
  \item 点Pと点Sの距離（地球儀上を移動する最短距離）は何cmか？
  \item 点Qと点Rの距離（地球儀上を移動する最短距離）は何cmか？
\end{enumerate}

\newpage %-------------------------------------------------- 段の区切り

%--------------------------------------------------------------- 6
\mondai $x$ ， $y$ が整数のとき，方程式 $7x-5y=1\ldots\text{①}$ を次の順序で解け。　（10点）

\vspace{3mm}

\begin{enumerate}
  \item ①を満たす1組の $(x,y)$ を求めよ。
  \item ①の一般解を求めよ。
\end{enumerate}

\vspace{18mm}

%--------------------------------------------------------------- 7
\mondai $29$ で割ると $20$ 余り， $37$ で割ると $21$ 余る自然数で，最小のものを求めよ。　（10点）

\vspace{18mm}

%--------------------------------------------------------------- 8
\mondai $x^2-3xy+2y^2=12$ を満たす自然数 $x$ ， $y$ の組をすべて求めよ。　（10点）

\vspace{18mm}

%--------------------------------------------------------------- 9
\mondai $\displaystyle\frac{1}{p}+\frac{1}{q}+\frac{1}{r}=1$ を満たす自然数 $p,q,r$ （ $p\leqq q\leqq r$ ）の組をすべて求めよ。　（10点）

\vspace{18mm}

%--------------------------------------------------------------- 10
\mondai $a+b+c+d=abcd$ を満たす自然数 $a,b,c,d$ の組は何組あるか。　（10点）

\vspace{18mm}

\end{document}
